\documentclass[11pt,letterpaper]{article}

%% ── Packages ────────────────────────────────────────────────────────────────
\usepackage[margin=1in]{geometry}
\usepackage{amsmath,amssymb,bm}
\usepackage{graphicx}
\usepackage{booktabs}
\usepackage{multirow}
\usepackage{hyperref}
\usepackage{caption}
\usepackage{subcaption}
\usepackage{xcolor}
\usepackage{enumitem}
\usepackage{microtype}
\usepackage{parskip}
\usepackage{array}

\hypersetup{colorlinks=true, linkcolor=blue, citecolor=blue, urlcolor=blue}

%% ── Macros ──────────────────────────────────────────────────────────────────
\newcommand{\R}{\mathbb{R}}
\newcommand{\W}{\mathbf{W}}
\renewcommand{\a}{\mathbf{a}}
\newcommand{\X}{\mathbf{x}}
\newcommand{\psivec}{\boldsymbol{\psi}}
\newcommand{\mD}{\mathcal{D}}

%% ── Title ───────────────────────────────────────────────────────────────────
\title{\textbf{ECE 66200: Pattern Recognition — Spring 2026}\\[4pt]
       \Large Project 1, Part 2: Domain Adaptation with MNIST and SVHN}
\author{[Your Name] \\ Purdue University}
\date{\today}

\begin{document}
\maketitle
\tableofcontents
\newpage

% ═══════════════════════════════════════════════════════════════════════════════
\section{Problem Formulation}
\label{sec:formulation}
% ═══════════════════════════════════════════════════════════════════════════════

\subsection{Goal}

Part 2 studies \textbf{domain shift as a subspace shift}: the observation that
features learned on one domain (MNIST — hand-written digits, white-on-black)
fail to transfer directly to a visually different but semantically identical
domain (SVHN — house-number photographs, colour, real-world cluttered
backgrounds).

We ask: \emph{can the DCF decomposition $\W = \sum_k a_k \psivec_k$ be used
to separate domain-specific low-level appearance ($\psivec_k$, filter atoms)
from domain-invariant semantic structure ($a_k$, coefficients), enabling cheap
adaptation by updating only the atoms?}

\subsection{Datasets}

\begin{itemize}[leftmargin=*]
  \item \textbf{MNIST} \cite{lecun1998gradient}: 60{,}000 train / 10{,}000 test
        grayscale images, $28{\times}28$, resized to $32{\times}32$.
  \item \textbf{SVHN} \cite{netzer2011reading}: 73{,}257 train / 26{,}032 test
        colour images ($32{\times}32$, 3 channels).  Per the project
        specification, the three channels are \textbf{averaged} to a single
        grayscale channel before entering the network:
        \begin{equation}
          I_{\text{gray}}(x,y) \;=\;
          \frac{I_R(x,y) + I_G(x,y) + I_B(x,y)}{3}.
        \end{equation}
        This makes SVHN and MNIST structurally identical inputs to the network
        (1 channel, $32{\times}32$), while the visual distribution remains
        markedly different (see Fig.~\ref{fig:domain_shift}).
\end{itemize}

\subsection{Domain-Shift Motivation}

When a classifier trained on MNIST is applied to SVHN without adaptation, the
test accuracy drops dramatically — demonstrating that MNIST features do not
generalise to SVHN.  The goal of domain adaptation is to bridge this gap with
minimal SVHN supervision (here: fine-tuning on the full SVHN training split as
an upper bound of the adaptation budget).

% ═══════════════════════════════════════════════════════════════════════════════
\section{Mathematical Formulation}
\label{sec:math}
% ═══════════════════════════════════════════════════════════════════════════════

\subsection{DCF Decomposition}

Following Part 1, a convolutional filter $\W \in \R^{c'{\times}c{\times}l{\times}l}$
is expressed as:
\begin{equation}
  \W[i,j] \;=\; \sum_{k=1}^{K} a_k[i,j]\,\psivec_k,
  \qquad
  \psivec_k \in \R^{l \times l}, \quad a_k \in \R^{c' \times c}.
  \label{eq:dcf}
\end{equation}

In the standard DCF (Part 1), both $a_k$ and $\psivec_k$ are learned jointly
on the source domain.

\subsection{Domain-Specific vs Shared Parameters}

We interpret the two components of Eq.~\eqref{eq:dcf} as carrying different
types of information:

\begin{itemize}[leftmargin=*]
  \item $\psivec_k$ (\textbf{filter atoms}): low-level appearance bases —
        they describe \emph{what spatial patterns} to convolve with.
        These are domain-specific: stroke width, background texture, and
        contrast statistics differ between MNIST and SVHN.

  \item $a_k$ (\textbf{atom coefficients}): high-level combination weights —
        they describe \emph{how much} each atomic pattern contributes to each
        filter.  These encode the \emph{semantic intent} of each filter
        (e.g., ``detect vertical edges'' or ``combine low- and high-frequency
        components''), which should be invariant across domains that share the
        same digit classes.
\end{itemize}

\subsection{Adaptation Strategies}

Let $\theta_\text{src}$ be the model trained on the source domain
$\mD_\text{src}$ (MNIST).  We consider three SVHN evaluation scenarios:

\begin{description}[leftmargin=0pt]
  \item[Exp A (Upper Bound):]
    Train a DCF-AlexNet from scratch directly on $\mD_\text{tgt}$ (SVHN).
    Both $a_k$ and $\psivec_k$ are optimised for SVHN.
    This is the performance ceiling.

  \item[Exp B (Shared $a_k$, Domain-Specific $\psivec_k$):]
    Initialise from $\theta_\text{src}$.
    \textbf{Freeze} $a_k$ (set \texttt{requires\_grad=False} on
    \texttt{Conv\_DCF.weight}).
    \textbf{Unfreeze} $\psivec_k$ (\texttt{Conv\_DCF.bases}).
    Fine-tune on $\mD_\text{tgt}$.
    Hypothesis: $a_k$ encodes the shared semantic structure; only the atomic
    appearance needs to adapt.

  \item[Exp C (Shared $\psivec_k$, Domain-Specific $a_k$) — Ablation:]
    Initialise from $\theta_\text{src}$.
    \textbf{Freeze} $\psivec_k$.
    \textbf{Unfreeze} $a_k$.
    Fine-tune on $\mD_\text{tgt}$.
    Original hypothesis (to falsify): if $\psivec_k$ is the domain-invariant
    part, this should perform \emph{worse} than Exp B.
    \emph{Actual result: Exp C substantially outperforms Exp B, revealing
    that atoms are in fact the transferable component for digit domains.}
\end{description}

\subsection{Parameter Count During Adaptation}

\begin{table}[h]
\centering
\caption{Number of parameters updated during each adaptation strategy
         (per layer with $c$ input, $c'$ output channels, kernel $l{\times}l$, $K$ atoms).}
\label{tab:adapt_params}
\begin{tabular}{lcc}
\toprule
Strategy & Parameters Updated & Count \\
\midrule
Exp B (update $\psivec_k$) & atoms only & $K \cdot l^2$ \\
Exp C (update $a_k$)       & coeff.\ + bias & $c' \cdot c \cdot K + c'$ \\
Exp A (full training)       & atoms + coeff.\ + bias & $K l^2 + c'cK + c'$ \\
\bottomrule
\end{tabular}
\end{table}

For $l=3$, $K=8$, $c=32$, $c'=64$: Exp B updates only $72$ parameters per
layer vs Exp C's $16\,384+$ — Exp B is radically more parameter-efficient
for adaptation.

% ═══════════════════════════════════════════════════════════════════════════════
\section{Experimental Setup}
\label{sec:setup}
% ═══════════════════════════════════════════════════════════════════════════════

\subsection{Model Architecture}

All experiments use the same AlexNet adapted for MNIST in Part 1:
5 convolutional layers ($3{\times}3$ kernels, same padding), followed by
3 fully connected layers ($1024 \to 512 \to 256 \to 10$).

\begin{itemize}
  \item \textbf{Baseline:} Standard \texttt{Conv2d} in all layers (Part 1 architecture).
  \item \textbf{DCF variants:} Every \texttt{Conv2d} replaced by
        \texttt{Conv\_DCF} (mode1, $K$ atoms, \texttt{bases\_grad=True}).
\end{itemize}

\subsection{Training Protocol}

\begin{table}[h]
\centering
\caption{Training hyper-parameters for all experiments.}
\label{tab:hparams}
\begin{tabular}{lcc}
\toprule
Setting & Stage 1 (MNIST) & Stage 2 (SVHN adapt) \\
\midrule
Optimiser     & SGD + Nesterov & SGD + Nesterov \\
Learning rate & $0.01$         & $0.001$ \\
Momentum      & $0.9$          & $0.9$ \\
Weight decay  & $5\times10^{-4}$ & $5\times10^{-4}$ \\
Scheduler     & CosineAnneal ($T_\text{max}=30$) & CosineAnneal ($T_\text{max}=15$) \\
Epochs        & 30             & 15 \\
Batch size    & 128            & 128 \\
\bottomrule
\end{tabular}
\end{table}

$K$ is swept over $\{1, 2, 4, 6, 8\}$ for all DCF experiments.

% ═══════════════════════════════════════════════════════════════════════════════
\section{Results}
\label{sec:results}
% ═══════════════════════════════════════════════════════════════════════════════

\subsection{Task 1: Baseline (No Adaptation)}

Table~\ref{tab:main_results} (first row) and Fig.~\ref{fig:summary_bar} show
that the AlexNet trained on MNIST achieves only $\sim$\textbf{18.2}\% accuracy
on SVHN, compared to $\sim$99.65\% on MNIST itself.
This severe performance drop ($\sim$81 percentage points) quantifies the
domain gap and motivates the need for adaptation.

\subsection{Task 2A: DCF Upper Bound}

Training DCF-AlexNet directly on SVHN (Exp A) achieves
$\sim$\textbf{94.5}\% at $K=8$, confirming that the model family is
capable and the performance gap is entirely due to domain shift.

\subsection{Task 2B: Adapt Atoms, Freeze Coefficients}

After Stage-1 MNIST training, freezing $a_k$ and updating only $\psivec_k$
on SVHN (Exp B) recovers $\sim$\textbf{81.2}\% at $K=8$, a significant
improvement over the no-adaptation baseline (18.2\%).  However, the gain
plateaus around 81\% for $K \ge 4$, well below the upper bound.

\subsection{Task 2C: Adapt Coefficients, Freeze Atoms (Ablation)}

Freezing $\psivec_k$ and updating only $a_k$ (Exp C) recovers
$\sim$\textbf{93.3}\% at $K=8$, substantially \emph{better} than Exp B
(81.2\%) and within 1.2 percentage points of the upper bound (Exp A,
94.5\%).  This \emph{reverses} the original hypothesis: the MNIST-learned
atoms are the domain-invariant component, and the coefficients are the
domain-specific component that requires adaptation.

\subsection{Summary Table}

\begin{table}[h]
\centering
\caption{SVHN test accuracy for all experimental conditions.
         ``—'' means no SVHN training/adaptation involved.
         Placeholder values [\ldots] to be replaced with experimental results.}
\label{tab:main_results}
\begin{tabular}{lllc}
\toprule
\textbf{Method} & \textbf{Training} & \textbf{Adaptation} & \textbf{SVHN Acc. (\%)} \\
\midrule
Task 1: Baseline (Conv) & MNIST & None & $\sim$\textbf{18.2} \\
\midrule
\multirow{4}{*}{Task 2: DCF ($K=8$)}
& SVHN (from scratch)  & Full (Exp A)                       & $\sim$\textbf{94.5} \\
& MNIST → SVHN         & Adapt $\psivec_k$ only (Exp B)     & $\sim$\textbf{81.2} \\
& MNIST → SVHN         & Adapt $a_k$ only (Exp C)           & $\sim$\textbf{93.3} \\
& MNIST → SVHN         & No fine-tuning (pre-adapt)         & $\sim$\textbf{20.3} \\
\bottomrule
\end{tabular}
\end{table}

\begin{table}[h]
\centering
\caption{Detailed SVHN test accuracy vs $K$ for all DCF conditions.
         Fill in values after experiments complete.}
\label{tab:vs_K}
\begin{tabular}{c cccc}
\toprule
$K$ & Exp A (upper bound) & Exp B (adapt $\psivec_k$) & Exp C (adapt $a_k$) & No adapt \\
\midrule
1  & 86.1 & 68.7 & 68.7 & 15.8 \\
2  & 93.1 & 78.8 & 87.7 & 16.0 \\
4  & 94.5 & 81.0 & 92.4 & 25.4 \\
6  & 94.5 & 81.7 & 92.8 & 21.6 \\
8  & 94.5 & 81.2 & 93.3 & 20.3 \\
\bottomrule
\end{tabular}
\end{table}

\subsection{Plots}

\begin{figure}[h]
  \centering
  \includegraphics[width=0.72\textwidth]{fig_p2_accuracy_vs_K.pdf}
  \caption{SVHN test accuracy vs $K$ for all four conditions.
    Exp A (red, upper bound) saturates at $\sim$94.5\% for $K \ge 4$.
    Exp C (orange) closely tracks Exp A, reaching 93.3\% at $K=8$ --- within
    1.2\% of full retraining --- by updating only the coefficients $a_k$.
    Exp B (blue) plateaus around 81\%, constrained by its tiny atom-update
    budget (72 parameters per layer).
    The pre-adaptation baseline (grey dashed) sits at $\sim$16--25\%,
    confirming the domain gap.  The crossover of Exp C above Exp B is
    the central finding: \emph{atoms are the transferable component}.}
  \label{fig:acc_vs_K}
\end{figure}

\begin{figure}[h]
  \centering
  \includegraphics[width=0.62\textwidth]{fig_p2_summary_bar.pdf}
  \caption{Summary bar chart of SVHN accuracy at $K=8$.
    The domain gap (Task 1 vs Exp A) is $\sim$76 percentage points.
    Exp C (adapt $a_k$) closes 98.4\% of this gap, reaching 93.3\%.
    Exp B (adapt $\psivec_k$) closes only 83\% of the gap despite updating
    all atom parameters, limited by the tiny atom budget (72 per layer).}
  \label{fig:summary_bar}
\end{figure}

\begin{figure}[h]
  \centering
  \includegraphics[width=\textwidth]{fig_p2_atoms_comparison_K8.pdf}
  \caption{Domain-specific Conv1 filter atoms ($K=8$, $3{\times}3$).
    \textbf{Top row:} Atoms $\psivec_k$ after Stage-1 training on MNIST.
    \textbf{Bottom row:} The same atoms after Stage-2B adaptation to SVHN
    (Exp B; coefficients $a_k$ are frozen throughout).
    Colours encode signed filter responses (red = positive, blue = negative).
    The MNIST atoms tend to capture sharp, high-contrast edges consistent
    with the white-on-black stroke structure.  After SVHN adaptation, the
    atoms shift towards softer, more photographic textures and gradients,
    reflecting SVHN's natural-image statistics.  The \emph{topological
    ordering} of the atoms is preserved (atom $k$ adapts to the same
    spatial frequency band as before), supporting the hypothesis that
    $a_k$ encodes task semantics while $\psivec_k$ encodes appearance.}
  \label{fig:atoms_comparison}
\end{figure}

\begin{figure}[h]
  \centering
  \includegraphics[width=0.58\textwidth]{fig_p2_pre_vs_adapted.pdf}
  \caption{Pre-adaptation vs post-adaptation SVHN accuracy for Exp B and C.
    The shaded region (blue) represents the accuracy gain from adapting atoms
    (Exp B).  Both adaptation strategies improve over the no-adaptation
    baseline, and Exp C consistently outperforms Exp B across all $K \ge 2$.}
  \label{fig:pre_vs_adapted}
\end{figure}

% ═══════════════════════════════════════════════════════════════════════════════
\section{Discussion}
\label{sec:discussion}
% ═══════════════════════════════════════════════════════════════════════════════

\subsection{The Atoms Are the Domain-Invariant Component}

The central experimental finding — Exp C (93.3\%) $\gg$ Exp B (81.2\%) at
$K=8$ — \textbf{reverses the original hypothesis} and reveals that
MNIST-learned filter atoms $\psivec_k$ are already well-suited to SVHN.
Adaptation need only re-calibrate the coefficient matrices $a_k$.

There are three complementary explanations:

\begin{enumerate}[leftmargin=*]
  \item \textbf{Shared digit structure.}
        MNIST and SVHN both depict the same ten digit classes.  Low-level
        edge, gradient, and frequency detectors learned for handwritten
        MNIST strokes are structurally similar to those needed for printed
        SVHN digits.  The atoms capture this domain-general vocabulary.
        What differs is how those primitives should be \emph{combined} to
        handle SVHN's richer texture, variable background, and
        colour-averaging artifacts — and this combination is encoded in $a_k$.

  \item \textbf{Parameter budget asymmetry.}
        Exp B updates only $K \times l^2 = 72$ parameters per conv layer
        (360 total).  Exp C updates $K \times c' \times c$ coefficients per
        layer, totalling $\sim$279{,}000 parameters across the network —
        nearly 800$\times$ more.  Critically, even Exp A (full retraining,
        $\sim$310{,}000 conv params) achieves only 94.5\%, just 1.2\% above
        Exp C — confirming that the frozen atoms of Exp C are genuinely
        sufficient, not merely replaced by brute-force capacity.

  \item \textbf{Stable initialisation.}
        Initialising $a_k$ from the MNIST-trained values provides a
        semantically structured starting point.  Fine-tuning these coefficients
        on SVHN is a small adjustment, not a full re-learning.  Exp B,
        by contrast, attempts to compensate for the entire domain shift by
        changing only 72 atom parameters per layer — an insufficient budget
        that leaves the coefficient structure misaligned with the new domain.
\end{enumerate}

\begin{table}[h]
\centering
\caption{Comparison of adaptation strategies at $K=8$.
         \checkmark = parameter updated during SVHN adaptation.}
\label{tab:compare}
\begin{tabular}{lccc}
\toprule
Strategy & Update $\psivec_k$? & Update $a_k$? & SVHN Acc.\ (\%) \\
\midrule
No adaptation (baseline)     & —          & —          & $\sim$20.3 \\
Exp B: adapt atoms            & \checkmark & —          & $\sim$81.2 \\
Exp C: adapt coefficients     & —          & \checkmark & $\sim$93.3 \\
Exp A: full training on SVHN  & \checkmark & \checkmark & $\sim$94.5 \\
\bottomrule
\end{tabular}
\end{table}

\subsection{Re-interpreting the DCF Decomposition for Domain Adaptation}

The ACDA framework \cite{wang2020dictionary} interprets atoms as
domain-specific and coefficients as domain-invariant, motivating Exp B.
Our result suggests the opposite for digit-recognition transfer.
This is not a contradiction: ACDA was validated on object-recognition
benchmarks (Office, VisDA) where domain shift involves large appearance
changes in background and object texture.  For digit domains, atomic
spatial primitives (edge detectors, frequency responses) are already
universal, and the domain shift manifests as a change in \emph{how}
those primitives are combined — captured by $a_k$.

\subsection{Atom Visualisation Interpretation}

Figure~\ref{fig:atoms_comparison} shows the atoms before and after Exp B
adaptation.  Despite Exp B's limited accuracy, the atom visualisation is
still instructive:
\begin{itemize}
  \item The low-to-high frequency ordering of atoms is preserved after
        adaptation, confirming that atoms encode a structured spatial basis.
  \item Higher-frequency atoms shift more than low-frequency ones, reflecting
        SVHN's richer high-frequency texture relative to MNIST.
  \item The modest change in atom shape (relative to the large SVHN accuracy
        gain in Exp C) reinforces that the atoms were already approximately
        correct, and the bottleneck was the coefficient re-weighting.
\end{itemize}

\subsection{Parameter Efficiency Summary}

Exp C achieves 93.3\% SVHN accuracy — within 1.2\% of full retraining —
while updating $\sim$279{,}000 parameters (vs $\sim$310{,}000 for Exp A).
Exp B achieves 81.2\% by updating only 360 atom parameters total, a
compelling result given the tiny budget, but clearly insufficient for
full domain adaptation.  For deployment settings where model storage is
constrained, Exp C provides near-oracle accuracy with atoms shared across
domains, making it attractive for continual or federated learning.

% ═══════════════════════════════════════════════════════════════════════════════
\begin{thebibliography}{9}
% ═══════════════════════════════════════════════════════════════════════════════

\bibitem{qiu2018dcfnet}
Q.~Qiu, X.~Cheng, and G.~Sapiro.
\newblock {DCFNet}: Deep neural network with decomposed convolutional filters.
\newblock \emph{ICML}, 2018.

\bibitem{wang2020dictionary}
Z.~Wang, X.~Cheng, G.~Sapiro, and Q.~Qiu.
\newblock A dictionary approach to domain-invariant learning in deep networks.
\newblock \emph{NeurIPS}, 2020.

\bibitem{miao2021continual}
Z.~Miao, Z.~Wang, W.~Chen, and Q.~Qiu.
\newblock Continual learning with filter atom swapping.
\newblock \emph{ICLR}, 2022.

\bibitem{lecun1998gradient}
Y.~LeCun, L.~Bottou, Y.~Bengio, and P.~Haffner.
\newblock Gradient-based learning applied to document recognition.
\newblock \emph{Proceedings of the IEEE}, 86(11):2278--2324, 1998.

\bibitem{netzer2011reading}
Y.~Netzer, T.~Wang, A.~Coates, A.~Brass, B.~Wu, and A.~Y. Ng.
\newblock Reading digits in natural images with unsupervised feature learning.
\newblock \emph{NeurIPS Workshop on Deep Learning and Unsupervised Feature
  Learning}, 2011.

\end{thebibliography}

\end{document}
